% Replace the current Chapter 2 with this version.
% Add references to the official SIAGE and INSIGHT project presentations before
% the final submission.

\chapter{General Context of the Project}

\section{Host institution and research setting}

This internship was carried out within the International Chair ``Inclusive
Societies and Ageing'' (SIAGE) at the University of Lorraine. The Chair studies
population ageing through the question of inclusive societies: how can social,
institutional, physical, and digital environments support the participation,
autonomy, rights, and citizenship of older adults? This perspective does not
reduce ageing to an individual health issue. It also considers life courses,
social inequalities, public policies, professional practices, and territorial
configurations that can facilitate or hinder inclusion.

SIAGE develops an interdisciplinary research environment that brings together
social sciences, public-policy studies, public health, engineering, and digital
technologies. Its objective is to connect academic research with territorial
actors, public decision-makers, professionals, associations, and older citizens.
Within this setting, digital methods are useful when they help organise and
compare dispersed knowledge, but their interpretation remains dependent on the
social and territorial context in which initiatives are developed.

The internship also contributes to the INSIGHT project. INSIGHT promotes the
connection of Natural Language Processing and Artificial Intelligence with
research in the Humanities and Social Sciences. \textit{Index\_Insight} is a
concrete application of this orientation: it uses NLP methods to make a large
and heterogeneous body of documentation on territorial adaptation to ageing more
accessible for research and future operational use.

\section{Senior Citizen Design and territorial knowledge}

One of the SIAGE Chair's participatory action-research activities is the
\textit{Senior Citizen Design} approach. Developed in France since 2018 in
collaboration with local authorities and the REIACTIS network, it seeks to
involve older adults as citizens rather than only as users or beneficiaries of
public action. The approach is structured around four connected stages:
participatory diagnosis, citizen-centred co-design, dialogue with public
decision-makers, and follow-up and evaluation.

The Senior Citizen Design Workshops produce action sheets in which participants
describe local issues, proposed responses, target publics, actors, and resources.
These documents are valuable because they preserve knowledge produced through
participatory work. However, they represent only one part of the relevant
documentary landscape. Numerous other initiatives are carried by municipalities,
associations, social-protection organisations, care providers, and networks
specialised in ageing. Their websites, reports, project descriptions, and
practical sheets also document locally situated responses to ageing-related
challenges.

\section{The role of Index\_Insight}

The aim of \textit{Index\_Insight} is not simply to create a collection of
documents. It is to capitalise territorial initiatives, make them visible and
comparable, and facilitate the circulation of experience between territories. A
structured index can support several uses: researchers can explore how local
ageing-related issues are described; territorial actors can identify relevant
examples; and participants in Senior Citizen Design Workshops can draw on
initiatives developed elsewhere when discussing local needs and possible
responses.

This objective creates a specific NLP problem. The relevant information is not
held in a homogeneous database: it is dispersed across sources with different
formats, vocabularies, levels of detail, and editorial structures. Titles,
territories, target audiences, initiative holders, dates, themes, and descriptions
may be explicit, implicit, absent, or located in different parts of a source
document. The task is therefore to transform these heterogeneous resources into
structured, traceable metadata without erasing their provenance or presenting
automatic inferences as established facts.

The design of the Index was informed by dialogue between its technical development
and the SIAGE research team. This dialogue contributes to the selection of sources,
the definition of metadata fields, the vocabulary used for quality control, and the
interpretation of the outputs. It is essential because the technical schema must
remain useful for the concrete questions addressed in participatory diagnosis,
co-design, dialogue with decision-makers, and the capitalisation of experience.

\section{Corpus, stakeholders, and constraints}

The corpus combines two complementary forms of documentation. The first consists
of action sheets produced by older adults themselves within Senior Citizen Design
Workshops implemented by SIAGE. The second consists of initiative descriptions
published by institutional, associative, and territorial actors involved in the
adaptation of territories to population ageing in France. The external sources were
selected in dialogue with the SIAGE team according to their relevance to the scope
of the Index and the availability of sufficiently substantial material for
structuring.

The intended users are correspondingly diverse: researchers studying ageing and
local public action, SIAGE participants and facilitators, local elected
representatives and practitioners, associations, and, in the longer term, older
citizens. Their needs are not limited to finding a document. They include comparing
initiative types, identifying territorial scales and stakeholders, locating examples
related to a particular issue, and understanding the limits of the available
documentation.

The main technical constraints arise from the organisation of the field itself.
The French gerontological landscape includes actors operating at national,
regional, departmental, intermunicipal, and local levels, with complementary
objectives in prevention, housing, mobility, health, participation, digital
inclusion, and social connection. This multi-level governance produces a naturally
fragmented documentary landscape. Relevant information may be distributed across
several pages of the same website, embedded in PDF documents, or expressed through
source-specific terminology. The collection, normalisation, and structured
extraction pipeline described in the following chapters is a response to this
documentary and institutional heterogeneity.
